\section{Selección de predictores y relevancia física}

La comparación sistemática de las once combinaciones de predictores evidenció que la
combinación 4, conformada por la temperatura superficial del mar (SST), la
precipitación global (tp) y los índices de la Oscilación Madden-Julian
(RMM1, RMM2) es la más adecuada para el pronóstico de precipitación
mensual en la región de Bocacosta.

Este resultado es coherente con la dinámica atmosférica regional descrita en el marco
teórico. La SST del Pacífico Oriental modula directamente la intensidad convectiva sobre
la Bocacosta mediante su influencia en la Zona de Convergencia Intertropical (ZCIT) y en
los vientos alisios \cite{dai2000global}. La precipitación global (tp), por su
parte, encapsula la organización espacial de los sistemas convectivos a escala sinóptica y
planetaria, actuando como un predictor de la actividad convectiva de gran
escala que antecede a los eventos locales. La inclusión de los índices de la MJO aporta
información sobre la variabilidad intraestacional organizada, la cual modula de forma
significativa la distribución temporal de las lluvias en el sur de Guatemala.

En contraste, las combinaciones que incorporaron variables adicionales como el
geopotencial (z), la humedad relativa y específica (r, q) y
las componentes del viento (u, v) no mejoraron el desempeño en gran manera, y en
varios casos lo redujeron. Esto sugiere que, si bien estas variables tienen respaldo
físico como predictores de la convección \cite{houze2014}, su inclusión en el modelo
LSTM pudo introducir colinealidad o ruido estadístico que dificultó el aprendizaje de
los patrones relevantes, incrementando el riesgo de sobreajuste. De igual manera, la
temperatura a 2 metros (2t) y la presión superficial (sp) no aportaron precisión
significativamente, lo que indica que la información climática que contienen ya está
suficientemente representada en la SST y la precipitación global para la región de
interés.

Este hallazgo tiene implicaciones prácticas relevantes: un modelo operativo con solo
tres tipos de predictores de reanálisis resulta computacionalmente eficiente,
reproducible y de fácil actualización para una institución como el INSIVUMEH.

\section{Desempeño global del modelo en el año 2025}

Las estadísticas globales obtenidas al pronosticar todos los meses del año 2025 muestran
un desempeño general satisfactorio. La correlación de Pearson de $r = 0.9166$ y la de
Spearman de $\rho = 0.9122$, ambas con $p \approx 0$, indican que el modelo captura
correctamente la fase y la tendencia de la variabilidad de la precipitación a lo largo
del año. El coeficiente de Nash-Sutcliffe de $\text{NSE} = 0.7440$ confirma que la
capacidad predictiva del modelo es superior a la de la media
climatológica como predictor de referencia, situándose en un rango considerado aceptable
para aplicaciones hidrológicas. El índice de Kling-Gupta de
$\text{KGE} = 0.7503$ refuerza esta conclusión, indicando que el modelo logra un
equilibrio razonable entre correlación, variabilidad y sesgo.

No obstante, el modelo exhibe una sobreestimación sistemática positiva de
$47.22$ mm, equivalente a un porcentaje de sesgo de
$20.17\ \%$. Esta tendencia a sobreestimar la precipitación puede atribuirse a que la transformación 
aplicada al objetivo, atenúa los valores extremos durante el entrenamiento,
induciendo al modelo a prever magnitudes moderadamente superiores a las observadas al
invertir las predicciones al espacio original en milímetros. El RMSE de 102.55\,mm
refleja principalmente el impacto de los meses con mayor precipitación acumulada
(mayo--octubre), donde los errores absolutos son naturalmente mayores dado el amplio
rango dinámico de la variable.

\section{Análisis del desempeño por mes}

El análisis mensual revela una marcada heterogeneidad en el desempeño del modelo a
lo largo del ciclo anual, que puede interpretarse en el contexto de la estacionalidad
climática de la Bocacosta.

\subsection*{Meses secos (diciembre--marzo)}

Durante la estación seca el modelo presenta errores absolutos bajos
(RMSE de 11--55\,mm), consistentes con los volúmenes de precipitación reducidos del
período. Sin embargo, las métricas de eficiencia como el KGE muestran valores negativos
o cercanos a cero en enero ($\text{KGE} = -0.269$), febrero ($\text{KGE} = -1.155$)
y diciembre ($\text{KGE} = -0.514$), así como correlaciones débiles o negativas. Esta
aparente contradicción se explica debido a que, cuando la precipitación observada es cercana a
cero, pequeñas desviaciones absolutas producen métricas de eficiencia relativa muy
desfavorables. En febrero y diciembre la correlación de Pearson es negativa
($r = -0.033$ y $r = -0.246$, respectivamente), lo que indica que el modelo no logra
reproducir la variabilidad espacial de la precipitación residual en estos meses. Esto
puede deberse a que los escasos eventos de lluvia en la estación seca están gobernados
por mecanismos locales o frontales de corta duración difíciles de capturar mediante
predictores de gran escala.

\subsection*{Meses de transición y comienzo de la estación lluviosa (abril--junio)}

El mes de abril muestra el mayor error relativo del año ($\text{RMSE} = 118.7$\,mm,
$\text{KGE} = 0.011$), con un sesgo positivo de 108.9\,mm. Esto corresponde al inicio
de la estación lluviosa, un período de alta variabilidad interanual en el que la fecha
del primer evento lluvioso significativo puede diferir sustancialmente entre años. A pesar de
esto la similitud entre el patrón observado y el predicho es uno de las más marcadas de todo
el año.  La SST y la MJO, si bien modulan la probabilidad de convección, no determinan con precisión
el momento exacto de la transición, lo que induce al modelo a anticipar precipitaciones
abundantes antes de que estas se produzcan. Mayo presenta un patrón similar
($\text{RMSE} = 134.1$\,mm, sesgo $= 96.6$\,mm), aunque con una correlación aceptable
de $r = 0.588$, indicando que el modelo captura la tendencia general pero sobreestima
las magnitudes.

Junio y julio representan los meses con mejor desempeño dentro de la estación lluviosa
(KGE de 0.653 y 0.649, respectivamente), con sesgos pequeños ($-17.9$ y $-2.0$\,mm).
Esto indica que el modelo ha aprendido adecuadamente los patrones de precipitación
asociados al inicio consolidado del monzón y al paso sistemático de
ondas del este, forzantes que se expresan con claridad en la señal de la SST y los
índices de la MJO.

\subsection*{Pico de la estación lluviosa (agosto--octubre)}

Agosto presenta el mejor desempeño individual de todos los meses lluviosos
($r = 0.848$, $\text{KGE} = 0.703$), coincidiendo con el segundo máximo del ciclo
bimodal de precipitación de la Bocacosta, asociado al pico de la ZCIT en el Pacífico
Oriental. Octubre muestra también una correlación satisfactoria ($r = 0.722$), aunque
el sesgo positivo elevado ($114.2$\,mm) sugiere que el modelo sobreestima la extensión
de los eventos extremos en este mes.

El mes de septiembre constituye el mayor desafío para el modelo ($\text{RMSE} =
222.7$\,mm, sesgo $= 211.1$\,mm, $\text{KGE} = 0.350$). Esta degradación puede
atribuirse a que septiembre coincide con el pico de la temporada de ciclones tropicales
en el Atlántico y el Pacífico Oriental, y con la actividad máxima de la ZCIT
\cite{herrera2017}.
Los predictores utilizados codifican correctamente la predisposición atmosférica a la
convección, pero no discriminan con suficiente precisión la ocurrencia e
intensidad de eventos extremos individuales dentro del mes, lo cual representa una
limitación intrínseca del enfoque de predicción mensual con predictores de gran escala.

\subsection*{Fin de la estación lluviosa (noviembre)}

Noviembre presenta una correlación espacial alta ($r = 0.891$), lo que indica que el
modelo captura bien la variabilidad entre píxeles. Sin embargo, el $\text{KGE} = 0.205$
y el sesgo de 64.2\,mm reflejan una sobreestimación sistemática durante la transición
hacia la estación seca. En este período, la señal de la SST y la MJO puede aún indicar
condiciones favorables para la convección mientras la precipitación observada decrece
por el avance de la subsidencia y la reducción de la humedad superficial. Al final,
se presenta una subestimación sistemática durante el mes de diciembre. Detectando aunque
en mayor medida la reducción de la precipitación al final del año.

\section{Análisis del ciclo anual medio}

El ciclo anual medio muestra que el modelo reproduce con buena fidelidad
la forma general de la curva de precipitación mensual climatológica de la Bocacosta: el
ascenso de mayo a junio, la canícula de julio-agosto, el segundo máximo en septiembre y
el descenso sostenido hacia diciembre. El modelo tiende, sin embargo, a suavizar el
primer pico (junio) y a amplificar el segundo (septiembre), lo cual es consistente con
el sesgo positivo detectado y con la dificultad para representar adecuadamente los
eventos extremos de septiembre.

La sobreestimación generalizada en los meses lluviosos se refleja en la curva
pronosticada, que se desplaza hacia arriba respecto a la observada durante casi todo el
año.

\section{Interpretación de los mapas espaciales}

La comparación mensual de los mapas de precipitación acumulada pronosticados por la
LSTM con los registrados por ENACTS revela que el modelo reproduce correctamente los
patrones espaciales dominantes de la Bocacosta: la concentración de precipitación sobre
la franja paralela a la cadena volcánica, el gradiente norte-sur desde la llanura
costera hacia el altiplano, y la extensión diferencial entre la estación seca y la
lluviosa. La arquitectura multi-torre con modulación por coordenadas geográficas
contribuye a esta especificidad espacial, permitiendo que el modelo aprenda sesgos
locales sistemáticos mediante el sesgo por píxel aprendible.

Sin embargo, en los meses de mayor dificultad (septiembre, octubre) la LSTM distribuye
la precipitación de manera más uniforme sobre toda la región, sin capturar los núcleos
de precipitación extrema que los datos ENACTS registran en puntos específicos del
territorio. Esto indica que la resolución del reanálisis ERA5, aunque superior a la de
modelos globales más gruesos, introduce un límite en la capacidad del modelo para
reproducir la variabilidad orográfica a escala de detalle propia de la Bocacosta.

\begin{comment}
\subsection{Limitaciones del modelo y perspectivas de mejora}

Con base en el análisis precedente, se identifican las siguientes limitaciones
principales del modelo desarrollado:

\begin{enumerate}
    \item \textbf{Sesgo positivo sistemático}: La sobreestimación del $20.17\,\%$
    requiere la implementación de una corrección de sesgo empírica o cuantílica antes
    de su aplicación operativa.

    \item \textbf{Dificultad en meses de transición estacional}: El modelo no captura
    con precisión el inicio y el fin de la estación lluviosa, lo que limita su utilidad
    para pronósticos de comienzo de temporada. La incorporación de índices de humedad
    del suelo o de anomalías de viento en 850\,hPa como predictores adicionales podría
    mejorar este aspecto.

    \item \textbf{Subestimación de la estructura extrema en septiembre}: Los eventos de
    precipitación extrema asociados a ciclones tropicales o a la actividad máxima de la
    ZCIT no son bien representados. El uso de predictores específicos de actividad
    ciclónica, o de índices de vórtice relativo, podría complementar la información de
    la SST y la MJO en este mes crítico.

    \item \textbf{Resolución espacial del reanálisis}: La cuadrícula ERA5 no resuelve
    completamente la orografía escarpada de la Bocacosta, introduciendo un límite en la
    habilidad predictiva del modelo para precipitación orográfica de detalle.

    \item \textbf{Período de evaluación limitado}: La validación se realizó sobre un
    único año (2025). Una evaluación multianual que abarque fases contrastantes del
    ENSO permitiría caracterizar de manera más robusta la habilidad del modelo bajo
    distintas condiciones climáticas de fondo.
\end{enumerate}

A pesar de estas limitaciones, los resultados demuestran que la arquitectura LSTM
multi-torre con predictores físicamente motivados constituye una herramienta viable y
complementaria a los modelos numéricos de predicción del tiempo para el pronóstico de
precipitación mensual en la región de Bocacosta, alcanzando niveles de correlación y
eficiencia que justifican su desarrollo hacia un sistema operativo de apoyo al INSIVUMEH.
\end{comment}
